\chapter{Analiza pe ferestre temporale}

Analiza prezentata in capitolele anterioare estimeaza o singura retea financiara
folosind toate observatiile disponibile. Aceasta abordare \textit{statica} are
avantajul utilizarii eficiente a datelor, dar ascunde potential modificarile
structurale ale pietei de-a lungul timpului. Piata financiara americana a traversat in esantionul analizat mai multe faze distincte, caracterizate prin niveluri diferite
de volatilitate si grade diferite de interconectare sectoriale.

In acest capitol este prezentata o extensie \textit{dinamica} a analizei, bazata
pe ferestre temporale rulante. Aceasta permite investigarea stabilitatii structurii
retelei si identificarea perioadelor in care dependentele conditionale se modifica
semnificativ.

\section{Metodologia ferestrelor temporale}

\subsection{Definitia ferestrei rulante}

Fie $w$ dimensiunea ferestrei (in zile tranzactionate) si $s$ pasul de deplasare.
Definim o succesiune de subperioade:
\begin{equation}
    \mathcal{I}_k = \{t_k, t_k + 1, \ldots, t_k + w - 1\}, \quad k = 1, 2, \ldots, K,
    \label{eq:window}
\end{equation}
unde $t_k = t_1 + (k-1) \cdot s$, iar $K = \lfloor (T - w)/s \rfloor + 1$ este
numarul total de ferestre. Pentru fiecare fereastra $\mathcal{I}_k$ se repeta
intreaga procedura:
\begin{enumerate}
    \item Se estimeaza un model VAR(1) pe observatiile din $\mathcal{I}_k$.
    \item Se calculeaza matricea de covarianta $\hat{\Sigma}^{(k)}$ a reziduurilor.
    \item Se estimeaza matricea de precizie $\hat{\Theta}^{(k)}$ prin GLasso,
    folosind aceeasi valoare a parametrului de regularizare ca in analiza
    principala, pentru a pastra comparabilitatea intre ferestre.
    \item Se construieste reteaua $G^{(k)} = (V, E^{(k)}, w^{(k)})$.
    \item Se calculeaza indicatori de centralitate, cu accent pe
    \textit{strength\_abs}.
\end{enumerate}

Aceasta procedura este implementata in scriptul \texttt{06\_time\_window\_network.R},
iar sinteza rezultatelor este realizata in \texttt{07\_write\_time\_window\_summary\_md.R}.

\subsection{Alegerea dimensiunii ferestrei}

Alegerea lui $w$ presupune un compromis fundamental:
\begin{itemize}
    \item \textit{Ferestre mici} (de exemplu, $w = 120$ zile $\approx$ 6 luni):
    permit surprinderea schimbarilor rapide ale structurii pietei, dar reduc
    numarul de observatii disponibile per fereastra, crescand incertitudinea
    estimatiei.
    \item \textit{Ferestre mari} (de exemplu, $w = 500$ zile $\approx$ 2 ani):
    ofera estimatii mai stabile, dar pot obscura schimbarile structurale rapide.
\end{itemize}

In proiect, fereastra a fost aleasa astfel incat sa mentina un raport
$w/n \gg 1$ (unde $n = 11$ este numarul de sectoare), conditie necesara pentru
estimarea stabila a modelului VAR si a matricei de covarianta.

\section{Indicatori dinamici ai retelei}

\subsection{Evolutia centralitatii sectoriale}

Fisierul \texttt{top\_sector\_per\_window.csv} sintetizeaza, pentru fiecare
fereastra, sectorul cu cea mai ridicata valoare a masurii
\textit{strength\_abs}. Aceasta informatie permite investigarea stabilitatii
sau variabilitatii pozitiei centrale in retea.

Daca acelasi sector apare frecvent in rolul de nod central, atunci pozitia sa
poate fi interpretata ca relativ stabila in esantionul analizat. Daca, in
schimb, identitatea nodului dominant se modifica intre ferestre, atunci acest
lucru sugereaza ca structura dependentei conditionale este sensibila la contextul
de piata si la subperioada analizata.

Din punct de vedere metodologic, aceasta analiza completeaza reteaua estimata pe
intregul esantion si arata ca centralitatea structurala poate fi studiata nu doar
static, ci si dinamic.


\section{Interpretarea dinamica a rezultatelor}

Analiza pe ferestre temporale adauga o dimensiune esentiala interpretarii
structurii retelei. Ea arata ca pozitiile centrale ale sectoarelor nu trebuie
privite ca proprietati complet fixe ale pietei, ci ca rezultate care pot varia
in timp.

Din punct de vedere statistic, aceasta observatie este importanta deoarece
sugereaza ca reteaua estimata pe intregul esantion reprezinta o agregare a unor
configuratii locale diferite. In consecinta, analiza dinamica nu contrazice
rezultatele statice, ci le completeaza, indicand in ce masura structura
dependentei conditionale ramane stabila sau se modifica intre subperioade.

\section{Limitele analizei dinamice}

Analiza pe ferestre temporale presupune ipoteza implicita ca structura retelei
este constanta in interiorul fiecarei ferestre. Aceasta ipoteza poate fi
nerealista in perioade de schimbari rapide. O extensie mai riguroasa ar fi
utilizarea unor modele cu parametri variabili continuu in timp (TVP-VAR),
care insa depasesc scopul acestei lucrari \cite{diebold2014}.

De asemenea, rezultatele analizei dinamice depind de alegerea lui $w$ si $s$.
Robustetea concluziilor ar trebui verificata pentru mai multe valori ale
acestor parametri, ceea ce constituie o directie naturala de dezvoltare ulterioara.

In ciuda acestor limite, analiza pe ferestre temporale este valoroasa pentru
ca transforma reteaua dintr-un obiect static intr-un obiect dinamic, mai fidel
naturii evolutive a pietelor financiare.

In plus, comparabilitatea intre ferestre depinde de mentinerea unei proceduri
de estimare coerente de la o fereastra la alta. Din acest motiv, utilizarea
aceleiasi penalizari in toate ferestrele are un rol metodologic important,
deoarece reduce variatia artificiala indusa de recalibrari diferite ale
parametrului de regularizare.